\begin{frame}
    \frametitle{Application to fusion systems}
    \begin{itemize}
        \item Proposed fusion energy facilities have a source rate high enough
            to activate materials, low enough to avoid significant composition
            change via depletion
        \item Activated nuclides decay and relase high energy photons after
            long after reactor shutdown
        \item Computing this dose is an important quantity for saftey and
            liscensing 
    \end{itemize}
    Transport-independent depletion was used in the workflow to
    calculate the shutdown dose rate in \cite{peterson_development_2024}
    
\end{frame}

\begin{frame}
    \frametitle{Application to fuel cycles}
    \begin{itemize}
        \item Depletion determines how loaded fuel compositin affect spent fuel
            compositions and related fuel cycle metrics.
        \item Global fuel cycle accounts for hundreds of reactors at once
            $\rightarrow$ impractical to run transport-coupled depletion to get
            spent fuel compositions
        \item A common approach is to use ``recipes'' that are based on
            depletion calculations for a specific reactor.
    \end{itemize}
    Transport-independent depletion was used in Cyclus
    \cite{huff_fundamental_2016}, an open source fuel cycle simulator, to
    provide real-time fuel depletion capabilities that is reactor-agnostic
    \cite{bachmann_open-source_2024}.
    
\end{frame}
